\documentclass[11pt]{article}
\usepackage[margin=1in]{geometry}
\usepackage{amsmath,amssymb}
\usepackage{booktabs}

\title{Laplacian Smoothing of Breaking Momentum (\texttt{SMOOTH\_BRK})}
\author{}
\date{}

\begin{document}
\maketitle

\section{Motivation}

Wave-breaking momentum $\mathcal{B}(x,y)$ in the WEC formulation is computed from
WaveWatch~III dissipation fields that are interpolated onto the ROMS grid. At coarse
resolution ($\Delta x \approx 30$\,m in L4PTSAL), the surfzone spans only 2--4 cells.
In complex 2D geometry---embayments and headlands---the setup pressure gradient cannot
balance the concentrated forcing, producing unrealistic nearshore currents.

Zeroing breaking momentum entirely (\texttt{BRK0}) removes the symptom but is
unphysical. Instead, we spread $\mathcal{B}$ over a physically motivated minimum
width using iterative Laplacian smoothing, representing the unresolved cross-shore
extent of roller propagation and turbulent bore momentum transfer.

Only $\mathcal{B}$ is smoothed, not the breaking dissipation $\varepsilon_b$.
The Akb mixing pathway is therefore unaffected.

\section{Discrete Laplacian smoother}

We apply the isotropic Laplacian stencil from Shchepetkin (used in KPP boundary
layer smoothing). At each iteration the field is updated as
\begin{equation}
  \mathcal{B}_{i,j} \;\leftarrow\; \mathcal{B}_{i,j}
  + c_1\bigl(F^x_{i+1,j} - F^x_{i,j} + F^\eta_{i,j+1} - F^\eta_{i,j}\bigr),
  \label{eq:update}
\end{equation}
where $c_1 = 3/16$, and $F^x$, $F^\eta$ are masked elementary differences with
an isotropy correction controlled by $c = 1/12$.

The elementary differences are
\begin{align}
  \widetilde{F}^x_{i,j}    &= \bigl(\mathcal{B}_{i,j} - \mathcal{B}_{i-1,j}\bigr)\,m^u_{i,j}, \\
  \widetilde{F}^\eta_{i,j} &= \bigl(\mathcal{B}_{i,j} - \mathcal{B}_{i,j-1}\bigr)\,m^v_{i,j},
\end{align}
where $m^u$, $m^v$ are the u- and v-point land masks (enforcing Neumann conditions
at the coastline; see \S\ref{sec:neumann}). The corrected fluxes are
\begin{align}
  F^\eta_{i,j} &= \widetilde{F}^\eta_{i,j}
    + c\bigl(\widetilde{F}^x_{i+1,j} + \widetilde{F}^x_{i,j-1}
            - \widetilde{F}^x_{i,j}   - \widetilde{F}^x_{i+1,j-1}\bigr), \\
  F^x_{i,j}   &= \widetilde{F}^x_{i,j}
    + c\bigl(\widetilde{F}^\eta_{i,j+1} + \widetilde{F}^\eta_{i-1,j}
            - \widetilde{F}^\eta_{i,j}   - \widetilde{F}^\eta_{i-1,j+1}\bigr).
\end{align}
The transversal correction terms restore isotropy: without them ($c=0$) the
stencil is a standard 5-point Laplacian with grid-aligned bias; with $c=1/12$
the effective weights are
\[
\frac{1}{32}\begin{pmatrix} 1 & 4 & 1 \\ 4 & 12 & 4 \\ 1 & 4 & 1 \end{pmatrix},
\]
which is the unique 9-point isotropic discrete Laplacian on a uniform grid.

\section{Conservation}

Summing (\ref{eq:update}) over all interior points, the right-hand side telescopes:
\begin{equation}
  \sum_{i,j} \bigl(F^x_{i+1,j} - F^x_{i,j}\bigr) = 0, \qquad
  \sum_{i,j} \bigl(F^\eta_{i,j+1} - F^\eta_{i,j}\bigr) = 0,
\end{equation}
because each flux appears once with positive sign and once with negative sign
in the interior sum. At domain boundaries, the Neumann condition
(zero normal derivative, \S\ref{sec:neumann}) ensures that boundary fluxes
vanish, so no momentum is created or destroyed:
\begin{equation}
  \sum_{i,j} \mathcal{B}_{i,j}^{(n+1)} = \sum_{i,j} \mathcal{B}_{i,j}^{(n)}.
\end{equation}
On a variable-metric grid, conservation holds to $O(\Delta x^2)$; on a
uniform grid it is exact.

\section{Neumann boundary conditions}
\label{sec:neumann}

At physical boundaries (coastline and open-boundary edges), we impose the
discrete Neumann condition
\begin{equation}
  \mathcal{B}_{\text{ghost}} = \mathcal{B}_{\text{interior}},
\end{equation}
i.e.\ the ghost-point value mirrors the adjacent interior value, giving a
zero normal gradient. This enters the smoother through two mechanisms:

\begin{enumerate}
\item \textbf{Masking.} Multiplying the elementary differences by $m^u_{i,j}$
  and $m^v_{i,j}$ zeroes any flux that crosses a land--ocean interface.
  Physically, momentum reflects off headlands rather than penetrating into
  land cells.

\item \textbf{Ghost-point fill.} At non-periodic physical boundaries
  (open boundaries, domain edges), ghost points are filled by copying the
  nearest interior value before each smoothing pass. At periodic boundaries,
  the MPI halo exchange fills the ghost points. At interior subdomain
  boundaries, \texttt{exchange\_xxx(brk)} synchronises ghost points between
  MPI ranks before each pass.
\end{enumerate}

Together, these ensure that the flux through every boundary face is zero,
which is the discrete analogue of $\partial\mathcal{B}/\partial n = 0$ and
is what guarantees conservation (\S3).

\section{Iterative application and spreading width}

The smoother is applied $N$ times in sequence, with an MPI exchange of
$\mathcal{B}$ between each pass. For a single pass on a uniform grid with
spacing $\Delta x$, the effective diffusion coefficient is
$\kappa = c_1\,\Delta x^2 = \tfrac{3}{16}\,\Delta x^2$, and the Gaussian
spreading width after $N$ passes is
\begin{equation}
  \sigma = \Delta x\,\sqrt{\frac{3N}{8}}, \qquad
  \text{FWHM} = 2\sqrt{2\ln 2}\;\sigma \approx 2.355\,\sigma.
\end{equation}
Representative values for $\Delta x = 30$\,m:

\medskip
\begin{center}
\begin{tabular}{@{}cccc@{}}
\toprule
$N$ & $\sigma$ (m) & FWHM (m) & MPI exchanges \\
\midrule
2  & 26 &  61 & 2 \\
4  & 37 &  87 & 4 \\
8  & 52 & 122 & 8 \\
\bottomrule
\end{tabular}
\end{center}

\medskip\noindent
The default is $N=4$ (FWHM $\approx 87$\,m, comparable to the surfzone width
at Duck94-scale bathymetry).

\section{Implementation summary}

\begin{itemize}
\item CPP flag \texttt{SMOOTH\_BRK} under \texttt{\#ifdef BREAKING}.
  Inactive when \texttt{BRK0} is defined.
\item Parameter \texttt{N\_brk\_smooth\,=\,4} (hardcoded in
  \texttt{wec\_forces\_tile}).
\item Smoothing is applied to $\mathcal{B}(i,j)$ at rho-points, after the
  \texttt{BRK0} block and before the projection to
  \texttt{brk2dx}/\texttt{brk2de}. Both the surface-stress path
  (\texttt{SURFACE\_BREAK}) and the body-force path (\texttt{BODY\_BREAK})
  benefit, since both read $\mathcal{B}(i,j)$.
\item Scratch arrays \texttt{FX\_brk}, \texttt{FE\_brk}, \texttt{FE1\_brk}
  are module-level allocatables, allocated in \texttt{init\_arrays\_wec\_tile}.
\end{itemize}

\end{document}
